\section{Evaluation Methodology}
\label{sec:method}

\subsection{The sec-triage environment}
Each episode is one alert on a web service. The hidden cause is one of eight explanations: three real attacks (credential stuffing, SQL-injection probing, DNS command-and-control), an insider exfiltrating data, an exposed admin panel, a known-vulnerable component, an authorized vulnerability scan, and a monitoring false positive. The hypothesis set adds \textit{other}. Ten read-only probes, costing 1 to 3 units each, cover authentication failures, HTTP patterns, source reputation, admin exposure, per-user egress, DNS characteristics, component inventory, change tickets, threat intelligence, and a generic runbook. The environment is designed so that shortcuts fail. Threat intelligence says only that an indicator is malicious, and three attacks share that answer. The runbook carries no information. Several benign causes look like attacks until the right probe is run. Each cause appears in three variants. In \textit{base}, outcomes are deterministic except for a 10\% threat-intelligence lag. \textit{Noise} adds 15\% transient 503 failures and a 30\% lag. In \textit{drift}, the incident is superseded after the second probe: the state version advances and the true cause switches. This gives 24 tasks, and each study runs every task three times per configuration. Each episode is served from its own loopback HTTP sandbox.

\subsection{Prediction tables}
The ledger and the selector need $P(o \mid h, a)$. We compare three sources. \emph{Designer} tables come from the environment's generating function and serve as an oracle reference. \emph{Empirical} tables are estimated by counting outcomes in LLM-free development episodes: an exhaustive script runs every probe on each (cause, variant) instance $k$ times, for $k \in \{1, 5, 20, 100\}$. It uses only base and noise variants, draws its seeds from a space disjoint from evaluation, and applies $\varepsilon = 0.01$ smoothing. A unit test enforces that the estimator never calls the generating function. \emph{Elicited} tables are produced by the planner LLM itself. Every table is written to disk, hashed, and frozen before any evaluation episode runs. The empirical tables tell us how many observations are enough, not how the method behaves in an unknown world: development and evaluation episodes come from the same generator.

\subsection{Configurations}
All configurations use the same model, prompt template, probes, budgets (10 tool calls, 10 cost units, 12 decisions, 900\,s), duplicate blocking, and verifier.
\begin{itemize}
  \item \textbf{ReAct-style}: the planner sees the raw history only and picks every probe.
  \item \textbf{Memory-only}: the planner also sees the ledger and picks every probe itself.
  \item \textbf{\hesp{}}: the controller picks the probe (EIG/cost or random), with or without showing its ranking, and with or without the finish guard or the controller-side stop.
\end{itemize}

\subsection{Models and serving}
Qwen2.5-7B-Instruct (bf16), Qwen2.5-32B-Instruct-AWQ, and Qwen2.5-72B-Instruct-AWQ in every study. From v0.8 on, the Llama 3.1 family is added: 8B-Instruct (bf16) and 70B-Instruct-AWQ-INT4. Models are served by vLLM on a single RTX~6000-class GPU at temperature 0.2, with a distinct seed per call.

\subsection{Pre-registration and statistics}
Each study's hypotheses, arms, primary endpoints, and analysis were written down and frozen, together with the source hash and the table hashes, before any of its episodes ran. The pipeline smoke tests run before freezing are disclosed and excluded from all analyses. The outcome measure is \emph{verified completion}, the fraction of episodes that end with a verified verdict. Every other ending counts as a failure. We report paired per-task differences with percentile task-cluster bootstrap intervals (2{,}000 resamples, 24 clusters). When a study has two primary endpoints, each gets a Bonferroni-adjusted 97.5\% interval. All secondary comparisons are exploratory and uncorrected. We also report descriptive security metrics. Two causes are benign: the authorized scan and the monitoring false positive. The other six require action. Among episodes that end with a verdict, the \emph{missed-attack} rate is the fraction of cases with an actionable true cause in which the verdict names a benign cause. The \emph{false-escalation} rate is the fraction of cases with a benign true cause in which the verdict names an actionable one. The \emph{wrong-cause} rate is the fraction of verdicts that name the wrong cause. The \emph{unresolved} rate is the fraction of all episodes without a verified verdict. Ground truth is the cause in effect at the moment of the verdict.
